\chapter{Methodology}
\label{sec:methodology}

% This chapter enumerates and discusses the specific steps and activities that will be performed by the proponent to accomplish the objectives of the research. The discussion covers the activities from pre-proposal to Final Thesis Writing. 

% Ethical issues surrounding the research and the steps to be taken by the research team to adhere to ethical guidelines and principles should be included when discussing each of the activities.

% Research activities include inquiry, survey, research, brainstorming, canvassing, consultation, review, interview, observe, experiment, design, test, document, and other similar tasks. The following sections present an example set of activities and methods for conducting computing research. \textbf{\textcolor{red}{However, note that your group's actual research activities and methods will be different depending on your intended research contributions. You should consult closely with your thesis adviser for proper guidance. The following sections need not appear in your actual thesis proposal.}}

% \section{Data Collection}

% This section, also called "Data Gathering", focus on activities that will enable the researchers to understand their users and their context better, or to analyze the operating environment where the proposed software may be deployed. The types of activities would include interviews, observations, surveys, focus group discussions, and review of secondary data (reports, videos, manuals) to enable the researchers to identify user requirements; learn the existing processes, strategies, rules and guidelines; and/or build the datasets and knowledge resources needed in the study.

% Ethical issues surrounding the collection of data from human participants, and from internet or 3rd party sources should be indicated. The steps to be taken by the proponents to ensure that ethical guidelines are followed should also be explained.

% \section{Software Design and Implementation}

% This section covers the design of data structures, knowledge bases, user interface, and algorithms. It includes implementation and related activities - tool selection, unit testing, integration testing, and function testing - that were covered in the course "Software Engineering".

% Your Adviser may ask you to separate the Design and the Implementation sections as needed. 

% You may also indicate a specific software development lifecycle model that you plan to adopt in developing your software, such as the Agile methodology.

% For machine learning types of research, this section may be replaced with "Model Training".

% \section{Validation}

% This section contains the different types of validation activities that you will conduct to evaluate the performance of your model or algorithm or software. It may include user perception studies, end user acceptance testing, usability test, and model validation.

% Ethical issues surrounding the collection of data from human participants should be indicated. The steps to be taken by the proponents to ensure that ethical guidelines are followed should also be explained.

% This section should include a discussion of the following if working with human participants:
% \begin{itemize}
%     \item Participant Selection. Who are your target participants? How will they be selected? What data will be collected?
%     \item Orientation. How will ethics be observed? What documents, e.g., validation protocol, will be shared with participants to orient them regarding the research and the validation procedure? Will there be an Informed Consent and Informed Assent? Are they found in the Appendix? Will there be any pre-tests or interviews to be administered prior to the actual experiment?
%     \item Experiment Proper (aka Procedure). How long is the experiment or how long should the user use your software? Is it one-on-one, or one-to-many? What will you do while the experiment is ongoing - observe the participants using a checklist, guide the participants, record the interaction/usage?
%     \item Debriefing. After the experiment, will you administer any post-tests or debriefing to solicit feedback from your participants? What are these?
% \end{itemize}

% \subsection{Calendar of Activities}

% For the proposal, a Gantt chart showing the schedule of the activities should be added. You can create a table like below, or create a figure. For example:

% Table \ref{tab:timetableactivities} shows a Gantt chart of the activities.  Each bullet represents approximately one week worth of activity.

% \begin{table}[]
%     \centering
%     \begin{tabular}{l|c|c|c|c|c|c|c|c|c}
%          \toprule
%          & \multicolumn{6}{c|}{2021} & \multicolumn{3}{c}{2022} \\
%          \textbf{Activity} & \textbf{Apr} & \textbf{May} & \textbf{Jun} & \textbf{Jul} & \textbf{Aug} & \textbf{Sep} & \textbf{Jan} & \textbf{Feb} & \textbf{Mar} \\
%          \midrule
%          Data Collection & **** & **** & & & & & & & \\
%          \midrule
%          Software Design & & **** & **** & **** & **** & **** & & & \\
%          \midrule
%          Validation & & & & & & **** & **** & **** & \\
%          \bottomrule
%     \end{tabular}
%     \caption{Gantt Chart of Activities}
%     \label{tab:timetableactivities}
% \end{table}

\hl{The proposed methodology follows a multi-layered framework designed to investigate the influence of content and context on short-form video engagement. While the current pipeline utilizes a specific Large Multimodal Model (LMM) architecture and a defined set of initial features, this approach is intended to be iterative and flexible. The research acknowledges that data representations and model configurations may be refined as the donated TikTok archives are processed. Consequently, while the procedures detailed in this section represent the primary technical path, the study maintains the scope to explore alternative modeling techniques or additional feature signals to ensure the most robust analysis of the engagement data.}

\begin{figure}[h]
    \centering
    \includegraphics[width=0.8\linewidth]{Revised Document/figures/chap3/methodology overview.png}
    \caption{Overview of the proposed methodology}
    \label{fig:methodology_overview}
\end{figure}

\section{Building the Dataset}

The construction of the dataset will address the limitations of content-only benchmarks like SnapUGC. These existing models often lack the creator-specific and contextual metadata required for cold-start prediction \cite{li2024delving}. Standard automated scraping typically results in inconsistent labels and restricted access to backend performance metrics \cite{li2025vquala}. This study will utilize a data donation approach where creators voluntarily provide first-party analytics to ensure the accuracy of the ground truth labels.

\subsection{Data Acquisition and Collection Procedure}

Target participants will consist of active TikTok content creators who maintain a professional or business account. \hl{The study aims to recruit an ideal minimum of 30 to 50 Filipino Micro-creators through a standardized recruitment process using a distributed Google Form link. Upon providing informed consent, participants will be required to donate their entire available video history (raw MP4 files) alongside their official TikTok analytics exports. This strategy ensures a high density of data points per user, allowing the model to analyze long-term creator patterns alongside specific video performance. The step-by-step acquisition process will follow a structured pipeline as summarized in } Figure \ref{fig:data_acqui}.

\hl{The data acquisition is executed through five distinct stages to ensure both technical integrity and participant privacy. The process begins with Informed Consent (Step 1), where creators review the research objectives and data handling policies before providing a digital signature. Once cleared, participants initiate an official Data Export (Step 2) from their TikTok account settings to generate a standardized archive containing their verified engagement logs.}

\hl{During the Submission (Step 3) phase, creators upload this archive and their raw MP4 files to a secure, encrypted cloud portal. To protect participant identity, researchers immediately perform Anonymization (Step 4), which involves stripping personal identifiers and replacing account handles with randomized alphanumeric IDs. The final stage is Verification (Step 5), where the timestamps from the analytics export are cross-referenced with the video metadata to ensure a precise mapping of engagement metrics to the specific multimodal content.}

\begin{figure}[h]
    \centering
    \includegraphics[width=1\linewidth]{Revised Document/figures/chap3/Data Acquisition.png}
    \caption{Step-by-Step Data Collection and Anonymization Procedure}
    \label{fig:data_acqui}
\end{figure}

% \begin{table}[H]
%     \centering
%     \begin{tabular}{lp{10cm}}
%         \toprule
%         \textbf{Step} & \textbf{Description} \\
%         \midrule
%         1. Consent & Participants will sign a digital agreement after reviewing the research objectives. \\
%         2. Export & Creators will generate an official "Download your data" archive from their TikTok settings. \\
%         3. Submission & Creators will upload the archive and raw MP4 files to a secure, encrypted cloud portal. \\
%         4. Anonymization & Researchers will strip all personal identifiers and replace creator handles with alphanumeric IDs. \\
%         5. Verification & Researchers will cross-reference timestamps to map engagement metrics to the video files. \\
%         \bottomrule
%     \end{tabular}
%     \caption{Step-by-Step Data Acquisition Procedure}
%     \label{tab:acquisition_procedure}
% \end{table} 

\subsection{Dataset Variables and Metrics}

The resulting dataset will integrate three primary feature layers \hl{, which serve as the initial framework for representing short-form video engagement.} The Multimodal Content Layer will consist of raw MP4 video files, captions, and hashtags. The Creator-Related Layer will incorporate metadata including follower counts and account age. The Contextual Layer will be represented by specific posting timestamps. This study will rely on the model to automatically extract and separate the visual and acoustic signals from the raw MP4 files. 

\hl{Acknowledging the potential for evolving data representations, this feature set remains flexible and subject to refinement based on the specific attributes available within the donated TikTok archives. While these represent the primary variables for the study, the methodology allows for the inclusion of additional signals.}
A comprehensive list of these \hl{initial} features is provided in Table \ref{tab:feature_list}.

\begin{table}[H]
    \centering
    \begin{tabular}{ll}
        \toprule
        \textbf{Feature}  & \textbf{Category} \\
        \midrule
        Raw MP4 File  & Multimodal Content \\
        Follower Count  & Creator-Related \\
        Creator Age  & Creator-Related \\
        Posting Time (Hour \& Day)  & Contextual \\
        \bottomrule
    \end{tabular}
    \caption{List of \hl{Initial} Dataset Features}
    \label{tab:feature_list}
\end{table}

\subsection{Ethical Safeguards and Limitations}

Ethical integrity will be maintained through strict de-identification protocols where all usernames, profile pictures, and handles will be removed immediately after collection. \hl{Regarding facial privacy and biometric data, the informed consent process will explicitly detail the processing of the creator’s likeness. To protect the privacy of non-consenting third parties appearing in the background of the videos, automated face-detection and blurring algorithms will be applied to secondary subjects to prevent unauthorized identification.}

TikTok Terms of Service compliance is ensured as the data collection is user-initiated through donation rather than automated scraping. To mitigate potential risks of privacy breaches or data misuse, all files will be stored in an encrypted environment. Furthermore, participants will retain the right to withdraw and request the deletion of their donated data at any time before the final model aggregation. While this approach provides granular data, it remains limited by a smaller sample size and potential voluntary response bias.

\section{Enhancing the Model}

The proposed study will implement an ensemble Large Multimodal Model (LMM) architecture to harmonize high-level semantic reasoning with structured creator metadata. This approach will transition from traditional two-stage regression networks to a unified framework capable of modeling complex cross-modality interactions.

\subsubsection{LMM Ensemble Architecture}
The framework will leverage an ensemble of state-of-the-art LMMs to capture diverse engagement signals:
\begin{itemize}
    \item \textbf{VideoLLaMA2 Branch:} An audio-visual-language model will be utilized for joint spatiotemporal and auditory understanding. It will process key video frames, background sound, and textual metadata simultaneously.
    \item \textbf{Qwen2.5-VL Branch:} A vision-language model will focus on general visual-semantic reasoning. It will provide complementary reasoning using keyframes and text-based metadata.
    \item \textbf{InternVideo2 Branch:} This branch will be integrated into the ensemble to ingest sampled keyframes for refined visual dynamics analysis.
\end{itemize}

\subsubsection{Integration of Creator and Contextual Layers}
A primary enhancement of this study will be the injection of ``non-content'' variables into the feature fusion layer to address the informational deficit in traditional cold-start modeling. While standard LMMs reason primarily about the media file, this model will augment those embeddings with structured metadata available at the time of upload:
\begin{itemize}
    \item \textbf{Creator Influence:} Follower count and account age will be integrated to account for the social influence and creator-audience relationships that bias initial distribution.
    \item \textbf{Contextual Timing:} Posting time will be incorporated to model the impact of peak user activity on early engagement probabilities.
\end{itemize}

\subsubsection{\hl{Technical Implementation Workflow}}
The execution of the model will follow a structured four-phase pipeline:

\paragraph{Phase 1: Preprocessing and Input Formatting}
The model will extract exactly eight keyframes from each video using uniform sampling to focus on the initial ``hook'' window. Visual inputs will be resized to a resolution of $384 \times 384$ for VideoLLaMA2, $768 \times 28 \times 28$ for Qwen2.5-VL and $224 \times $224 for InternVideo2 to ensure memory efficiency \cite{li2025vquala} . Background audio will be converted into spectrograms to capture emotional and narrative cues.

\paragraph{Phase 2: Feature Extraction and Backbone Freezing}
To manage computational requirements, the study will utilize the LMMs primarily as feature extractors. The vision and audio encoders of the VideoLLaMA2 branch will remain frozen during training to preserve pre-trained multimodal priors. The model will extract the hidden representations from the final layer of the LLM decoders to serve as high-dimensional content embeddings.

\paragraph{Phase 3: Creator-Oriented Feature Fusion}
The study will employ a fusion layer where high-dimensional content embeddings are concatenated with the scalar metadata of the creator. Creator features, such as follower count, will undergo log-transformation or Min-Max normalization to ensure they are on a comparable scale with the neural embeddings. A cross-modal attention mechanism will then merge these inputs into a unified video representation.

\paragraph{Phase 4: Joint Training and Regression}
The final prediction will be handled by a lightweight regression head composed of a two-layer multi-layer perceptron (MLP). This head will include a dropout layer and a fully connected layer with 2048 neurons. The model will be trained to jointly predict \textit{Engagement Continuation Rate} (ECR) and \textit{Normalized Average Watch Percentage} (NAWP). The joint training loss $L$ is defined as:

\begin{equation}
    L = ||NAWP - \widehat{NAWP}||_2 + ||ECR - \widehat{ECR}||_2
\end{equation}

This joint regression strategy will leverage the 0.926 correlation between initial attention and overall retention to stabilize the learning process and improve prediction accuracy.

% Enhancing the model part above


\section{Evaluating the Model}

To evaluate the performance of our enhanced model, it will be compared to current benchmark models. It will be evaluated using the same metrics as those models.

\subsection{Evaluation Metrics}

The model's objective is to predict the ECR and NAWP of a given short video. ECR is the probability that the user will watch a video after 5 seconds. NAWP measures the retention of the user, regardless of video length.

This study will assess the model using the Root Mean Squared Error (RMSE) of both ECR and NAWP, the Spearman Rank Correlation Coefficient (SRCC), and the Pearson Linear Correlation Coefficient (PLCC). SRCC assesses the ranking accuracy between the actual and predicted scores. PLCC measures the linear agreement and consistency between the actual and predicted scores. A final score is computed using both SRCC and PLCC: 
\begin{equation}
    \text{Final Score} = 0.6\times SRCC + 0.4\times PLCC
\end{equation}

This final score provides a balanced score of both ranking consistency and prediction accuracy.

\subsection{Benchmark Comparison}

Our study will use the models by \citeauthor{sun2025engagement} and \citeauthor{guan2025mmf} as our benchmarks for comparison. We chose these because they are the best performing models available in predicting user engagement for SFVs in cold-start scenarios \cite{li2025vquala}.

As shown in Figure \ref{fig:sun2025_framework}, \citeauthor{sun2025engagement} utilized an ensemble of two LMMs: VideoLLaMA2 \cite{cheng2024videollama}, which is an LMM designed for audio-visual and spatio-temporal understanding, and Qwen2.5-VL \cite{bai2025qwen3}, which is a vision-language model.

\begin{figure}[h]
    \centering
    \includegraphics[width=1\linewidth]{chap3/sun2025_framework.png}
    \caption[Sun et al., 2025 Framework Overview]
    {Framework Overview. Source: \cite{sun2025engagement}.}
    \label{fig:sun2025_framework}
\end{figure}

On the other hand, \citeauthor{guan2025mmf} predicted ECR with the use of a collaborative framework of three branches as shown in Figure \ref{fig:guan2025_framework}: (1) a foundational model extracting multimodal features, (2) a disentangled model for technical and aesthetic factors based on DOVER \cite{wu2023exploring}, and (3) an instruction-tuned Qwen2.5-VL model.

\begin{figure}[h]
    \centering
    \includegraphics[width=1\linewidth]{chap3/guan2025_framework.png}
    \caption[Guan et al., 2025 Framework Overview]
    {Framework Overview. Source: \cite{guan2025mmf}.}
    \label{fig:guan2025_framework}
\end{figure}

\subsection{Experimental Studies}

To determine the impact of the proposed features on prediction accuracy, ablation studies will be conducted. These experiments will isolate or combine the different categories from Table \ref{tab:feature_list}: multimodal content $F_{content}$, creator features $F_{creator}$, and contextual $F_{context}$, into various models $M_n$. In doing so, each factor's impact on the prediction accuracy of the ECR and NAWP is quantitatively measured. Table \ref{tab:experiment_configurations} shows the configurations of each experiment we will do.

\begin{table}[h]
    \centering
    \begin{tabular}{lccc}
        \toprule
        \textbf{Model} & $F_{content}$ & $F_{creator}$ & $F_{context}$ \\
        \midrule
         $M_0$: Baseline & \checkmark &  &  \\
         $M_1$: Creator-Augmented & \checkmark & \checkmark &  \\
         $M_2$: Context-Aware & \checkmark &  & \checkmark \\
         $M_3$: Fully Integrated & \checkmark & \checkmark & \checkmark \\
        \bottomrule
    \end{tabular}
    \caption{Experimental Configurations for Ablation Studies}
    \label{tab:experiment_configurations}
\end{table}

\hl{We can then apply statistical tests to assess whether the improvements between our model and the baseline models significantly enhance the prediction accuracy.}

\section{Calendar of Activities}

The proposed methodology will be done over the span of 12 months, divided into the three phases previously outlined.

The first phase will focus on gathering and preparing the data needed for the study. This includes selecting a group of creators, collecting their consent and data, and finally anonymizing, verifying, and cleaning the data. This will span from May 2026 to August 2026. 

The second phase will focus on feature extraction, model training, enhancement, and initial evaluation. Training on the benchmark models will also occur here. This will span from August 2026 to December 2026.

The last phase will be the evaluation of the proposed model and performing ablation studies. This will span from January 2027 to April 2027.

Table \ref{tab:gantt_chart} shows a Gantt Chart of the activities.


\begin{table}[h]
    \centering
    \resizebox{\linewidth}{!}{\begin{tabular}{l|c|c|c|c|c|c|c|c|c|c|c|c}
         \toprule
         & \multicolumn{8}{c|}{2026} & \multicolumn{4}{c}{2027} \\
         \textbf{Activity} & \textbf{May} & \textbf{Jun} & \textbf{Jul} & \textbf{Aug} & \textbf{Sep} & \textbf{Oct} & \textbf{Nov} & \textbf{Dec} & \textbf{Jan} & \textbf{Feb} & \textbf{Mar} & \textbf{Apr} \\
         \midrule
         Building the Dataset & **** & **** & **** & **** & & & & & & & & \\
         \midrule
         Enhancing the Model & & & & **** & **** & **** & **** & **** & & & & \\
         \midrule
         Evaluating the Model & & & & & & & & **** & **** & **** & **** & **** \\
         \bottomrule
    \end{tabular}}
    \caption{Gantt Chart of Activities}
    \label{tab:gantt_chart}
\end{table}